\chapter{Discussion}
\label{chap: disussion}

The preceding chapter demonstrates that the proposed hybrid pipeline reproduces the statistical observables of the $Re = 2500$ cylinder wake: the time-averaged velocity field, the three independent Reynolds-stress components, and the integrated force coefficients, all with appropriate physical accuracy relative to the reference simulation. On its own, these results show that latent neural ODE rollouts can replace part of the conventional time integration without compromising the quantities that matter most to the flow. They do not, however, say whether the replacement was worth it. This is what the acceleration metrics $S_\text{loop}$ and $S_\text{eff}$ aim to show. This chapter aims to provide context for the acceleration and accuracy results reported in the previous section, to specify the conditions under which the reported speedup actually holds, and to clarify where it does not. Several of those results look favourable or disfavourable only because certain choices were kept fixed, such as the grid size and the hardware on which the pipeline was run. These choices are examined and discussed in turn to adequately assess the results of this thesis.


\section{Interpreting the acceleration metrics}
\label{sect: disc_metrics}

The speedup metrics constructed in this work are key indicators of the current framework's functionality. The hybrid-loop speedup for the base simulation $S_\text{loop} = T_\text{wall}^\text{ref} /T_\text{wall}^\text{h, tot} \approx 1.47$ measures the cost of hybrid integration \emph{once the networks already exist}; the effective speedup $S_\text{eff} = T_\text{wall}^\text{ref} / (T_\text{wall}^\text{h, tot} + T_\text{wall}^\text{h, training}) \approx 0.07$ charges the workflow for building those networks. The gap between the two is the central pitfall of this work: for a single run over the integrated window, with two full retraining procedures (Table~\ref{tab: wall_timing_hybrid}; the third update is data-only), the complete pipeline is roughly an order of magnitude \emph{slower} than conventionally advancing the reference solver. The per-CTU rollout is some $84.6$ times cheaper than conventional time integration of a single CTU of CFD (\autoref{fig: base_cost_ctu}), but this advantage is dwarfed by the fixed cost of training: the autoencoder and neural ODE together consume $682$~s, equivalent to roughly $1018$~CTU of reference integration, against the $100$~CTU actually simulated.

The disparity between the two metrics is itself informative. The loop speedup $S_\text{loop} > 1$ confirms that, once trained, $\mathcal{P}$ integrates the flow more cheaply per CTU than the reference solver; the effective speedup $S_\text{eff} < 1$ shows that this per-step speedup is still overshadowed by the cost of network training. These two metrics together do show that $\mathcal{P}$ holds the capacity for genuine end-to-end acceleration, but the realisation of it depends on shifting the balance between these costs, either by reducing the training expense through smarter network training regimes and better reuse of the pipeline, or by raising the cost per CTU of the reference simulation so that each replaced step is worth more. 

Generalisation experiments act on the two factors that set $S_\text{loop}$ separately: the cost of each replaced step and the number of steps a rollout replaces. A finer grid raises the per-CTU cost of the memory-bound reference solver~\cite{Weymouth2025} while leaving the latent-rollout cost fixed, so each replaced step is worth more and the $S_\text{loop}$ ceiling rises. The change in Reynolds number changes the number of steps that get replaced by $\mathcal{P}$, but in an inverse way: lower $Re$ yields smoother, lower-dimensional dynamics that $f_\phi$ could in principle follow for longer, yet the realised rollout length is governed by the KNN gate rather than by $f_\phi$ alone. An increase in Reynolds number makes the latent dynamics more complex, making them harder for $f_\phi$ to learn and leading to increased drift. But the KNN monitor adopts a more permissive gate due to increased variance in the latent training set, which enables longer rollouts and thus better acceleration performance. 

The honest statement of the results of the developed framework is therefore conditional: the hybrid loop is faster than CFD per step, but the workflow as a whole would only serve as a true acceleration method for longer, more expensive or repeated simulation windows. This conditionality also has an upside. Because the rollout cost is fixed by the latent dimension while the cost of a CFD step grows with the simulation size, the per-ctu cost margin widens as the reference solver becomes more expensive, as demonstrated in the finer-mesh case of \autoref{subsect: 4_RE5000}. The regimes in which $\mathcal{P}$ is worth its training cost are therefore precisely the more demanding regimes, where the CFD step is costly enough that saving and replacing it repays the training overhead more quickly. Given that the KNN monitor is adequately fit to the training set, this is discussed further in \autoref{sect: disc_monitor}. 

\section{The role of hardware in the reported gain}
\label{sect: disc_hardware}

A second important discussion note is hardware-related. With the exception of the autoencoder training, every component in $\mathcal{P}$ (the WaterLily reference, the warmup and stabilisation CFD steps, the neural ODE training, and the latent rollout) is executed on CPU. This means that the $S_\text{loop}$ is a clean CPU-versus-CPU comparison on identical hardware, which is appropriate for isolating the algorithmic gain. However, since WaterLily is a backend-agnostic solver, it can run on GPU architectures, where its benchmarks report approximately 2 orders of magnitude of speedup on GPU relative to the serial CPU execution \cite{Weymouth2025}. 

If the reference had been run on a GPU, the per-CTU CFD cost would fall while the cost of a latent rollout would remain the same, eroding the performance of $S_\text{loop}$. This erosion is, however, strongly grid-dependent: the ${\sim}2$ orders of magnitude reported in \cite{Weymouth2025} are an asymptotic figure, reached only once the GPU memory is completely utilised at large grids, whereas the present $256^2$ case ($\,{\sim}0.07$~M DOF) lies below the smallest benchmarked grid and would see a considerably smaller GPU advantage. Even at this reduced factor, because the rollout cost is fixed, GPU execution of the reference would push $S_\text{loop}$ towards, and plausibly below, break-even on equal hardware. The GPU comparison should therefore be read as the most demanding case the hybrid loop can be asked to handle, not as a representative one.


Additionally, the accelerated workflow was deliberately kept on the CPU, mainly because training $f_\phi$ is not expected to benefit greatly from GPU acceleration: the network is small, and the cost is dominated by the sequential ODE solve and the reverse-mode differentiation through it, which parallelise poorly across the latent batch. This means that for GPU operation of $\mathcal{P}$, data would constantly need to be transferred between the CPU and GPU memory, adding an overhead to $\mathcal{P}$ that would spoil the maximal acceleration that it could possibly achieve. 


However, this does not invalidate the approach. WaterLily was selected as the development vehicle for this framework precisely because it has differentiable and backend-agnostic properties that enable the coupling described in Section~\ref{sect: coupling}. Many of the CFD solvers used in practice, however, are not yet able to utilise GPU architectures; for these, CPU execution would be the only path to acceleration, as described in this work. The GPU comparison above should therefore serve as the most demanding case the hybrid loop can be asked to face, rather than the representative one: in any setting where the reference solver is CPU-bound by construction, the per-step gain reported through $S_\text{loop}$ is the gain that genuinely applies, and WaterLily serves here only as a convenient and well-characterised stand-in for a broader class of solvers.


\section{Sensitivity across parameter space}
\label{sect: disc_params}

The developed framework exposes a large number of coupled hyperparameters, and the various runs in \autoref{app: ae_arch} varied only a small number of them around a single baseline. The autoencoder alone introduces a large number of parameters that could influence the eventual solution, more than those addressed in this work (kernel size, activation functions, optimisation algorithm, etc.). The addition of the neural ODE does not make the problem easier; it adds a whole new set of parameters that are coupled with those of the autoencoder. This coupling is the central difficulty. The latent geometry that $f_\phi$ must learn is itself produced by $\varphi_\theta$, so a change to the autoencoder (a wider bottleneck $N_z$, a different physics weighting $(\lambda_\text{div},\,\lambda_\text{curl})$, or simply a longer training budget) reshapes the very trajectories $\boldsymbol{\mathcal{T}}_z$ on which the multiple-shooting loss, the group size $n_s$, and the continuity weight $\lambda_c$ act. The KNN monitor sits one step further downstream of this: its threshold $\tau$ and neighbour count $k$ are calibrated on $\boldsymbol{\mathcal{T}}_z$, so any shift in the latent representation propagates through to the OOD decision and, with it, the rollout length that ultimately governs $S_\text{loop}$. The parameters of $\mathcal{P}$, therefore, cannot be regarded as independent dials; an optimum found for one component is conditional on the state of the others.

The sensitivity study of Section~\ref{sect: 4_p_parameters} respects this only partially. The one-factor-at-a-time (OFAT) procedure cannot reveal the complex interactions between the parameters of $\mathcal{P}$. A full accounting of parameter interactions would require sampling the entire hyperparameter space, and the number of pipeline trainings this demands grows exponentially with the parameter count, which places it beyond the computational scope of this work. The results reported here should accordingly be interpreted as a marginal sensitivity study around a single operating point, not as a converged optimum. Meaning that the framework should be understood as one whose accessible performance is, in all likelihood, understated because the true optimal performance is not yet found. 

A further consequence of this coupling is visible in the comparison metrics themselves. The configurations that produced the largest $S_\text{loop}$ were not the most accurate but the least: the smallest bottleneck ($N_z = 8$) and the most lightly trained autoencoder both extended the average rollout, because an under-resolved encoder yields a lower-complexity latent signal that $f_\phi$ is able to keep within the training distribution for longer before the KNN monitor flags it. The same configurations were the worst at recovering the second-order statistics, with the shear stress $\langle u'v'\rangle$ and the transverse normal stress $\langle v'v'\rangle$ degrading most. Acceleration and statistical fidelity, therefore, pull against one another, and $S_\text{loop}$ on its own is a misleading objective: maximising it rewards discarding precisely the turbulent information $\mathcal{P}$ wants to preserve. The appropriate target is the trade-off between the two, and the operating point on that curve is dictated by whether acceleration or statistical accuracy is the priority for a given application, rather than by either metric alone.

\section{Generalisation to the \texorpdfstring{$Re=5000$}{Re=5000} regime}
\label{sect: disc_re5000}

The $Re=5000$, $512^2$ case (\autoref{subsect: 4_RE5000}) tests $\mathcal{P}$ on a simultaneously more expensive and more demanding operating point. The two changes act on different parts of the framework but were applied together: the grid refinement raises the per-CTU cost of the memory-bound reference solver while leaving the rollout cost fixed, and the higher Reynolds number drives the wake from quasi-periodic shedding into a more complex state. 

On cost, the grid refinement behaves as the memory-bound argument predicts: the reference rises to \SI{4842}{\milli\second}/CTU against \SI{866}{\milli\second}/CTU at $256^2$ while the rollout stays fixed at ${\sim}\SI{48}{\milli\second}$/CTU, lifting $S_\text{loop}$ to $2.07$ and $S_\text{eff}$ to $\approx0.19$ against $0.073$ at the base point, so the economic claim of \autoref{sect: disc_metrics} strengthens here. On fidelity, the same run exposes the limit of the compression: the autoencoder no longer generalises beyond its training window (validation and test losses of $9.6\times10^{-3}$ and $8.9\times10^{-2}$, \autoref{fig: RE5000_loss}), and the reconstruction cannot reconstruct the complex, high-gradient, near-wall layer that get integrated to the lift forces (\autoref{fig: RE5000_forces}). The favourable $S_\text{loop}$ and the degraded forces are two faces of one run, where the reconstruction, not the rollout, becomes the binding error source.

Compounding the poor generalisation outside the training window, the encoder's latent dynamics mimic the chaotic nature of the wake. This makes the training of $f_\phi$ harder due to the absence of the clean, near-periodic shedding signal that the $Re=2500$ latent trajectory offered the network to lock onto. Where the quasi-periodic base wake gave $f_\phi$ a low-dimensional, repetitive target it could follow well beyond the training horizon, the $Re=5000$ latent trajectory is highly chaotic, so the compact network has no stable structure to extrapolate, and the rollout diverges from the reference sooner. This is where the regime begins to expose the limit of the neural ODE itself: a more chaotic flow yields latent dynamics whose predictability, rather than the width of the bottleneck encoding them, sets how far $f_\phi$ can be trusted to advance the state.


\section{Integration Monitor}
\label{sect: disc_monitor}

The generalisation to $Re=250$ and $Re=5000$ expose a structural flaw in the KNN gate that is invisible at the base point. The threshold $\tau$ is set as a quantile of the in-distribution nearest-neighbour distances of $\boldsymbol{\mathcal{T}}_z$, so it scales with how dispersed the training distribution is. The well-resolved $Re=250$ wake is cleanly periodic and low-variance: its latent cluster is tight, $\tau\approx0.07$, and the gate is so strict that every rollout is clipped just past the threshold and reset, holding the realised speedup below what $f_\phi$'s near-perfect fit could otherwise sustain (\autoref{subsect: 4_re250}). The badly-resolved $Re=5000$ wake does the opposite: its chaotic latent distribution inflates $\tau$ from $0.614$ to $0.902$ across the retraining stages, the gate turns lenient, and rollouts are allowed to run for thousands of steps, exactly where $f_\phi$ accumulates the most drift (\autoref{subsect: 4_RE5000}).

The integration monitor is therefore more permissive where the dynamics are hardest to learn and predict, and most restrictive where they are easiest, which is the wrong way around. Since the acceleration framework must preserve accuracy, it should penalise incorrect rollouts in the regimes where the dynamics are most complex, where prediction errors are made quickly. This is not a tuning error but a direct consequence of deriving the trigger from training-set variance: the harder the flow is to learn, the wider its latent distribution, and the more latitude the gate gives to the part of the pipeline least able to use it. The deeper issue is that the KNN score measures proximity to $\boldsymbol{\mathcal{T}}_z$ in latent space, not the physical correctness of the decoded field. A latent point can sit comfortably inside a broad cluster while decoding to a physically wrong field, as at $Re=5000$, and can drift marginally past a tight $\tau$ while decoding to a perfectly acceptable field, as at $Re=250$.


Among these generalisation cases, the base case occupies the middle ground. The $Re=2500$, $256^2$ latent distribution carries enough variance that $\tau$ gives the rollout a useful margin to drift, rather than clipping a well-fitted rollout the moment it leaves the training cloud as at $Re=250$; while the latent space stays well-resolved enough that a genuinely unseen wake state still falls outside the cluster and crosses the threshold before the decoded field becomes unphysical, unlike at $Re=5000$. The monitor behaves as intended here, but this is a property of the operating point, not the gate: the base-case variance happens to land in the narrow window where a threshold read off the training spread coincides with a physically sensible one.

Circling back to \ref{sect: KNN}, the choice of the KNN score was deliberate, and in hindsight double-edged. It was selected precisely because it imposes no distributional assumption on the latent states: since the encoder is retrained per flow and yields a different latent geometry each time, any assumed parametric distribution would not transfer across cases, whereas a non-parametric neighbour distance adapts to whatever cluster the training data happens to form~\cite{Sun2022}. This generalisability leads to inverse gating behaviour over $Re$ variation, because a threshold read off the latent trajectory's own spread inherits the spread's dependence on the flow. A monitor gated on a physical residual instead, such as the decoded-field divergence $\langle\,|\hat{u}_{i,i}|\,\rangle$, would not share this defect: it judges every rollout against the same physical tolerance regardless of latent geometry, imposing accuracy directly rather than by proxy. This route was not chosen in this work due to the cost of decoding and evaluating a residual in physical space at every step, which is almost as expensive as advancing WaterLily itself, which would erase the speedup the rollout exists to provide, so a practical physical monitor hinges on finding a residual proxy cheap enough to check often.


